%% SPDX-License-Identifier: CC-BY-4.0
%% Copyright (C) 2026  Alin-Adrian Anton <alin.anton@upt.ro>
%%
%% diode-quickstart.tex -- the short path to a working one-way link.
%% Build:  make quickstart   (or  latexmk -pdf diode-quickstart.tex)

\documentclass[11pt,a4paper]{article}
\newcommand{\projectname}{opc-diode}
\newcommand{\doctitleshort}{Data Diode Quick Start}
%% SPDX-License-Identifier: CC-BY-4.0
%% Copyright (C) 2026  Alin-Adrian Anton <alin.anton@upt.ro>
%%
%% Shared with the sibling diode projects.  The canonical copy lives in
%% log-diode/docs/tutorial/ -- fix it there and re-copy, so opc-diode,
%% gnat-lt-pro, wal-postgresqldb and wal-mysqldb stay in step.
%% common.tex -- shared preamble for the diode tutorial documents.
%% Input from the preamble of both diode-quickstart.tex and
%% data-diode-handbook.tex.  Everything here must work with a stock
%% TeX Live / texlive-collection-latexrecommended install; no external
%% images, all figures are TikZ.

\usepackage[T1]{fontenc}
\usepackage{lmodern}
\usepackage{microtype}
\usepackage[a4paper,margin=25mm,headheight=15pt]{geometry}
\usepackage{amsmath,amssymb}
\usepackage{xcolor}
\usepackage{graphicx}      % \reflectbox (\copyleft), \resizebox for wide figures
\usepackage{booktabs,tabularx,longtable,array}
\usepackage{enumitem}
\usepackage{caption}
\usepackage{fancyhdr}
\usepackage{titlesec}
\usepackage{listings}
\usepackage{tikz}
\usetikzlibrary{positioning,arrows.meta,shapes.geometric,calc,fit,backgrounds}
\usepackage{tcolorbox}
\tcbuselibrary{skins,breakable}
\usepackage[hidelinks,pdfusetitle]{hyperref}

%% ---------------------------------------------------------------- colours
\definecolor{ldblue}  {HTML}{1F4E79}
\definecolor{ldred}   {HTML}{A61B1B}
\definecolor{ldgreen} {HTML}{1E6F3C}
\definecolor{ldamber} {HTML}{9A6700}
\definecolor{ldgray}  {HTML}{5A5A5A}
\definecolor{codebg}  {HTML}{F6F6F4}

%% ---------------------------------------------------------------- headers
\pagestyle{fancy}
\fancyhf{}
\fancyhead[L]{\footnotesize\textsc{\doctitleshort}}
\fancyhead[R]{\footnotesize\thepage}
\fancyfoot[C]{\footnotesize \projectname{} diode tutorial --- CC BY 4.0}
\renewcommand{\headrulewidth}{0.4pt}
\renewcommand{\footrulewidth}{0pt}

%% ---------------------------------------------------------------- listings
\lstset{
  basicstyle=\ttfamily\footnotesize,
  backgroundcolor=\color{codebg},
  frame=single,
  rulecolor=\color{ldgray!40},
  framesep=4pt,
  xleftmargin=6pt,
  xrightmargin=2pt,
  breaklines=true,
  breakatwhitespace=false,
  postbreak=\mbox{\textcolor{ldgray}{$\hookrightarrow$}\space},
  showstringspaces=false,
  keepspaces=true,
  columns=fullflexible,
  commentstyle=\color{ldgreen}\itshape,
  keywordstyle=\color{ldblue}\bfseries,
  stringstyle=\color{ldred},
  captionpos=b,
  aboveskip=8pt,
  belowskip=6pt,
  tabsize=2,
}

\lstdefinestyle{shell}{
  language=bash,
  morekeywords={ip,nft,iptables,ip6tables,pfctl,sysctl,systemctl,tcpdump,arp,
                ethtool,openssl,hping3,arping,nc,socat},
}
\lstdefinestyle{conf}{
  language=,
  morecomment=[l]{\#},
}
\lstdefinestyle{nft}{
  language=,
  morecomment=[l]{\#},
  morekeywords={table,chain,rule,type,hook,priority,policy,iif,oif,iifname,
                oifname,udp,dport,daddr,saddr,accept,drop,filter,ingress,
                input,output,forward,inet,netdev,ip,ip6,meta,counter,limit,log},
  keywordstyle=\color{ldblue}\bfseries,
}
\lstdefinestyle{pf}{
  language=,
  morecomment=[l]{\#},
  morekeywords={set,block,pass,in,out,quick,on,proto,from,to,port,all,
                no-state,state,antispoof,skip,block-policy,scrub,keep},
  keywordstyle=\color{ldblue}\bfseries,
}
\lstdefinestyle{ada}{
  language=Ada,
  morekeywords={SPARK_Mode,Pre,Post,Loop_Invariant},
}
\lstdefinestyle{c}{
  language=C,
  basicstyle=\ttfamily\scriptsize,
  morekeywords={ssize_t,size_t,sig_atomic_t,uint8_t,IFF_TAP,IFF_NO_PI,
                TUNSETIFF,POLLIN,O_RDWR,O_CLOEXEC,EINTR,EAGAIN,IFNAMSIZ},
  numbers=left,
  numberstyle=\ttfamily\tiny\color{ldgray!70},
  numbersep=7pt,
}

%% ------------------------------------------------------------- admonitions
\newtcolorbox{warnbox}[1][Warning]{%
  breakable, enhanced, colback=ldred!4, colframe=ldred,
  boxrule=0.9pt, left=6pt, right=6pt, top=5pt, bottom=5pt,
  fonttitle=\bfseries\sffamily\small, title={#1}}

\newtcolorbox{labbox}[1][Laboratory use only]{%
  breakable, enhanced, colback=ldamber!6, colframe=ldamber,
  boxrule=0.9pt, left=6pt, right=6pt, top=5pt, bottom=5pt,
  fonttitle=\bfseries\sffamily\small, title={#1}}

\newtcolorbox{notebox}[1][Note]{%
  breakable, enhanced, colback=ldblue!4, colframe=ldblue,
  boxrule=0.7pt, left=6pt, right=6pt, top=5pt, bottom=5pt,
  fonttitle=\bfseries\sffamily\small, title={#1}}

\newtcolorbox{tipbox}[1][In practice]{%
  breakable, enhanced, colback=ldgreen!4, colframe=ldgreen,
  boxrule=0.7pt, left=6pt, right=6pt, top=5pt, bottom=5pt,
  fonttitle=\bfseries\sffamily\small, title={#1}}

%% ------------------------------------------------------------------ macros
%% Copyleft (U+1F12F).  No Latin Modern / T1 glyph exists and neither
%% cclicenses nor ccicons is in a base TeX Live, so mirror the copyright
%% sign -- which is exactly what the symbol is.  \mathord keeps the spacing
%% sane when it lands next to text.
\newcommand{\copyleft}{\reflectbox{\copyright}}

\newcommand{\cmd}[1]{\texttt{#1}}
\newcommand{\file}[1]{\texttt{#1}}
\newcommand{\prog}[1]{\textsf{#1}}
%% Renewed per project, immediately after %% SPDX-License-Identifier: CC-BY-4.0
%% Copyright (C) 2026  Alin-Adrian Anton <alin.anton@upt.ro>
%%
%% Shared with the sibling diode projects.  The canonical copy lives in
%% log-diode/docs/tutorial/ -- fix it there and re-copy, so opc-diode,
%% gnat-lt-pro, wal-postgresqldb and wal-mysqldb stay in step.
%% common.tex -- shared preamble for the diode tutorial documents.
%% Input from the preamble of both diode-quickstart.tex and
%% data-diode-handbook.tex.  Everything here must work with a stock
%% TeX Live / texlive-collection-latexrecommended install; no external
%% images, all figures are TikZ.

\usepackage[T1]{fontenc}
\usepackage{lmodern}
\usepackage{microtype}
\usepackage[a4paper,margin=25mm,headheight=15pt]{geometry}
\usepackage{amsmath,amssymb}
\usepackage{xcolor}
\usepackage{graphicx}      % \reflectbox (\copyleft), \resizebox for wide figures
\usepackage{booktabs,tabularx,longtable,array}
\usepackage{enumitem}
\usepackage{caption}
\usepackage{fancyhdr}
\usepackage{titlesec}
\usepackage{listings}
\usepackage{tikz}
\usetikzlibrary{positioning,arrows.meta,shapes.geometric,calc,fit,backgrounds}
\usepackage{tcolorbox}
\tcbuselibrary{skins,breakable}
\usepackage[hidelinks,pdfusetitle]{hyperref}

%% ---------------------------------------------------------------- colours
\definecolor{ldblue}  {HTML}{1F4E79}
\definecolor{ldred}   {HTML}{A61B1B}
\definecolor{ldgreen} {HTML}{1E6F3C}
\definecolor{ldamber} {HTML}{9A6700}
\definecolor{ldgray}  {HTML}{5A5A5A}
\definecolor{codebg}  {HTML}{F6F6F4}

%% ---------------------------------------------------------------- headers
\pagestyle{fancy}
\fancyhf{}
\fancyhead[L]{\footnotesize\textsc{\doctitleshort}}
\fancyhead[R]{\footnotesize\thepage}
\fancyfoot[C]{\footnotesize \projectname{} diode tutorial --- CC BY 4.0}
\renewcommand{\headrulewidth}{0.4pt}
\renewcommand{\footrulewidth}{0pt}

%% ---------------------------------------------------------------- listings
\lstset{
  basicstyle=\ttfamily\footnotesize,
  backgroundcolor=\color{codebg},
  frame=single,
  rulecolor=\color{ldgray!40},
  framesep=4pt,
  xleftmargin=6pt,
  xrightmargin=2pt,
  breaklines=true,
  breakatwhitespace=false,
  postbreak=\mbox{\textcolor{ldgray}{$\hookrightarrow$}\space},
  showstringspaces=false,
  keepspaces=true,
  columns=fullflexible,
  commentstyle=\color{ldgreen}\itshape,
  keywordstyle=\color{ldblue}\bfseries,
  stringstyle=\color{ldred},
  captionpos=b,
  aboveskip=8pt,
  belowskip=6pt,
  tabsize=2,
}

\lstdefinestyle{shell}{
  language=bash,
  morekeywords={ip,nft,iptables,ip6tables,pfctl,sysctl,systemctl,tcpdump,arp,
                ethtool,openssl,hping3,arping,nc,socat},
}
\lstdefinestyle{conf}{
  language=,
  morecomment=[l]{\#},
}
\lstdefinestyle{nft}{
  language=,
  morecomment=[l]{\#},
  morekeywords={table,chain,rule,type,hook,priority,policy,iif,oif,iifname,
                oifname,udp,dport,daddr,saddr,accept,drop,filter,ingress,
                input,output,forward,inet,netdev,ip,ip6,meta,counter,limit,log},
  keywordstyle=\color{ldblue}\bfseries,
}
\lstdefinestyle{pf}{
  language=,
  morecomment=[l]{\#},
  morekeywords={set,block,pass,in,out,quick,on,proto,from,to,port,all,
                no-state,state,antispoof,skip,block-policy,scrub,keep},
  keywordstyle=\color{ldblue}\bfseries,
}
\lstdefinestyle{ada}{
  language=Ada,
  morekeywords={SPARK_Mode,Pre,Post,Loop_Invariant},
}
\lstdefinestyle{c}{
  language=C,
  basicstyle=\ttfamily\scriptsize,
  morekeywords={ssize_t,size_t,sig_atomic_t,uint8_t,IFF_TAP,IFF_NO_PI,
                TUNSETIFF,POLLIN,O_RDWR,O_CLOEXEC,EINTR,EAGAIN,IFNAMSIZ},
  numbers=left,
  numberstyle=\ttfamily\tiny\color{ldgray!70},
  numbersep=7pt,
}

%% ------------------------------------------------------------- admonitions
\newtcolorbox{warnbox}[1][Warning]{%
  breakable, enhanced, colback=ldred!4, colframe=ldred,
  boxrule=0.9pt, left=6pt, right=6pt, top=5pt, bottom=5pt,
  fonttitle=\bfseries\sffamily\small, title={#1}}

\newtcolorbox{labbox}[1][Laboratory use only]{%
  breakable, enhanced, colback=ldamber!6, colframe=ldamber,
  boxrule=0.9pt, left=6pt, right=6pt, top=5pt, bottom=5pt,
  fonttitle=\bfseries\sffamily\small, title={#1}}

\newtcolorbox{notebox}[1][Note]{%
  breakable, enhanced, colback=ldblue!4, colframe=ldblue,
  boxrule=0.7pt, left=6pt, right=6pt, top=5pt, bottom=5pt,
  fonttitle=\bfseries\sffamily\small, title={#1}}

\newtcolorbox{tipbox}[1][In practice]{%
  breakable, enhanced, colback=ldgreen!4, colframe=ldgreen,
  boxrule=0.7pt, left=6pt, right=6pt, top=5pt, bottom=5pt,
  fonttitle=\bfseries\sffamily\small, title={#1}}

%% ------------------------------------------------------------------ macros
%% Copyleft (U+1F12F).  No Latin Modern / T1 glyph exists and neither
%% cclicenses nor ccicons is in a base TeX Live, so mirror the copyright
%% sign -- which is exactly what the symbol is.  \mathord keeps the spacing
%% sane when it lands next to text.
\newcommand{\copyleft}{\reflectbox{\copyright}}

\newcommand{\cmd}[1]{\texttt{#1}}
\newcommand{\file}[1]{\texttt{#1}}
\newcommand{\prog}[1]{\textsf{#1}}
%% Renewed per project, immediately after %% SPDX-License-Identifier: CC-BY-4.0
%% Copyright (C) 2026  Alin-Adrian Anton <alin.anton@upt.ro>
%%
%% Shared with the sibling diode projects.  The canonical copy lives in
%% log-diode/docs/tutorial/ -- fix it there and re-copy, so opc-diode,
%% gnat-lt-pro, wal-postgresqldb and wal-mysqldb stay in step.
%% common.tex -- shared preamble for the diode tutorial documents.
%% Input from the preamble of both diode-quickstart.tex and
%% data-diode-handbook.tex.  Everything here must work with a stock
%% TeX Live / texlive-collection-latexrecommended install; no external
%% images, all figures are TikZ.

\usepackage[T1]{fontenc}
\usepackage{lmodern}
\usepackage{microtype}
\usepackage[a4paper,margin=25mm,headheight=15pt]{geometry}
\usepackage{amsmath,amssymb}
\usepackage{xcolor}
\usepackage{graphicx}      % \reflectbox (\copyleft), \resizebox for wide figures
\usepackage{booktabs,tabularx,longtable,array}
\usepackage{enumitem}
\usepackage{caption}
\usepackage{fancyhdr}
\usepackage{titlesec}
\usepackage{listings}
\usepackage{tikz}
\usetikzlibrary{positioning,arrows.meta,shapes.geometric,calc,fit,backgrounds}
\usepackage{tcolorbox}
\tcbuselibrary{skins,breakable}
\usepackage[hidelinks,pdfusetitle]{hyperref}

%% ---------------------------------------------------------------- colours
\definecolor{ldblue}  {HTML}{1F4E79}
\definecolor{ldred}   {HTML}{A61B1B}
\definecolor{ldgreen} {HTML}{1E6F3C}
\definecolor{ldamber} {HTML}{9A6700}
\definecolor{ldgray}  {HTML}{5A5A5A}
\definecolor{codebg}  {HTML}{F6F6F4}

%% ---------------------------------------------------------------- headers
\pagestyle{fancy}
\fancyhf{}
\fancyhead[L]{\footnotesize\textsc{\doctitleshort}}
\fancyhead[R]{\footnotesize\thepage}
\fancyfoot[C]{\footnotesize \projectname{} diode tutorial --- CC BY 4.0}
\renewcommand{\headrulewidth}{0.4pt}
\renewcommand{\footrulewidth}{0pt}

%% ---------------------------------------------------------------- listings
\lstset{
  basicstyle=\ttfamily\footnotesize,
  backgroundcolor=\color{codebg},
  frame=single,
  rulecolor=\color{ldgray!40},
  framesep=4pt,
  xleftmargin=6pt,
  xrightmargin=2pt,
  breaklines=true,
  breakatwhitespace=false,
  postbreak=\mbox{\textcolor{ldgray}{$\hookrightarrow$}\space},
  showstringspaces=false,
  keepspaces=true,
  columns=fullflexible,
  commentstyle=\color{ldgreen}\itshape,
  keywordstyle=\color{ldblue}\bfseries,
  stringstyle=\color{ldred},
  captionpos=b,
  aboveskip=8pt,
  belowskip=6pt,
  tabsize=2,
}

\lstdefinestyle{shell}{
  language=bash,
  morekeywords={ip,nft,iptables,ip6tables,pfctl,sysctl,systemctl,tcpdump,arp,
                ethtool,openssl,hping3,arping,nc,socat},
}
\lstdefinestyle{conf}{
  language=,
  morecomment=[l]{\#},
}
\lstdefinestyle{nft}{
  language=,
  morecomment=[l]{\#},
  morekeywords={table,chain,rule,type,hook,priority,policy,iif,oif,iifname,
                oifname,udp,dport,daddr,saddr,accept,drop,filter,ingress,
                input,output,forward,inet,netdev,ip,ip6,meta,counter,limit,log},
  keywordstyle=\color{ldblue}\bfseries,
}
\lstdefinestyle{pf}{
  language=,
  morecomment=[l]{\#},
  morekeywords={set,block,pass,in,out,quick,on,proto,from,to,port,all,
                no-state,state,antispoof,skip,block-policy,scrub,keep},
  keywordstyle=\color{ldblue}\bfseries,
}
\lstdefinestyle{ada}{
  language=Ada,
  morekeywords={SPARK_Mode,Pre,Post,Loop_Invariant},
}
\lstdefinestyle{c}{
  language=C,
  basicstyle=\ttfamily\scriptsize,
  morekeywords={ssize_t,size_t,sig_atomic_t,uint8_t,IFF_TAP,IFF_NO_PI,
                TUNSETIFF,POLLIN,O_RDWR,O_CLOEXEC,EINTR,EAGAIN,IFNAMSIZ},
  numbers=left,
  numberstyle=\ttfamily\tiny\color{ldgray!70},
  numbersep=7pt,
}

%% ------------------------------------------------------------- admonitions
\newtcolorbox{warnbox}[1][Warning]{%
  breakable, enhanced, colback=ldred!4, colframe=ldred,
  boxrule=0.9pt, left=6pt, right=6pt, top=5pt, bottom=5pt,
  fonttitle=\bfseries\sffamily\small, title={#1}}

\newtcolorbox{labbox}[1][Laboratory use only]{%
  breakable, enhanced, colback=ldamber!6, colframe=ldamber,
  boxrule=0.9pt, left=6pt, right=6pt, top=5pt, bottom=5pt,
  fonttitle=\bfseries\sffamily\small, title={#1}}

\newtcolorbox{notebox}[1][Note]{%
  breakable, enhanced, colback=ldblue!4, colframe=ldblue,
  boxrule=0.7pt, left=6pt, right=6pt, top=5pt, bottom=5pt,
  fonttitle=\bfseries\sffamily\small, title={#1}}

\newtcolorbox{tipbox}[1][In practice]{%
  breakable, enhanced, colback=ldgreen!4, colframe=ldgreen,
  boxrule=0.7pt, left=6pt, right=6pt, top=5pt, bottom=5pt,
  fonttitle=\bfseries\sffamily\small, title={#1}}

%% ------------------------------------------------------------------ macros
%% Copyleft (U+1F12F).  No Latin Modern / T1 glyph exists and neither
%% cclicenses nor ccicons is in a base TeX Live, so mirror the copyright
%% sign -- which is exactly what the symbol is.  \mathord keeps the spacing
%% sane when it lands next to text.
\newcommand{\copyleft}{\reflectbox{\copyright}}

\newcommand{\cmd}[1]{\texttt{#1}}
\newcommand{\file}[1]{\texttt{#1}}
\newcommand{\prog}[1]{\textsf{#1}}
%% Renewed per project, immediately after \input{common}.
\newcommand{\amp}{\prog{sender}}
\newcommand{\deamp}{\prog{receiver}}
%% The two ends of the link, named by ROLE rather than by security level.
%% A diode can be deployed in either orientation -- ingesting into a
%% protected domain, or exporting out of one -- so the mechanics here are
%% written in terms of which end sends and which receives. Each project's
%% "Orientation" section maps these onto its own security levels.
\newcommand{\srcside}{\textsc{source}}
\newcommand{\dstside}{\textsc{destination}}
\newcommand{\nm}[1]{#1\,nm}
\newcommand{\dbm}[1]{#1\,dBm}
\newcommand{\dbb}[1]{#1\,dB}

%% ------------------------------------------------------------ tikz styles
\tikzset{
  host/.style     ={draw, rounded corners=2pt, minimum width=27mm,
                    minimum height=12mm, align=center, font=\small,
                    fill=ldblue!8, draw=ldblue!70},
  mc/.style       ={draw, minimum width=25mm, minimum height=13mm,
                    align=center, font=\scriptsize, fill=gray!8,
                    draw=ldgray},
  netbox/.style   ={draw, dashed, rounded corners=3pt, draw=ldgray!70,
                    inner sep=5mm},
  splitter/.style ={draw, diamond, aspect=1.5, inner sep=1.5pt,
                    align=center, font=\scriptsize, fill=ldgreen!12,
                    draw=ldgreen},
  fiber/.style    ={-{Latex[length=2.4mm,width=1.8mm]}, thick,
                    draw=ldgreen!65!black},
  copper/.style   ={-, thick, draw=ldblue!70!black},
  dataflow/.style ={-{Latex[length=3.4mm,width=2.6mm]}, line width=1.5pt,
                    draw=ldblue},
  blocked/.style  ={-{Latex[length=2.4mm]}, thick, draw=ldred, dashed},
  lbl/.style      ={font=\scriptsize\sffamily, inner sep=1.5pt},
  lblg/.style     ={font=\scriptsize\sffamily, inner sep=1.5pt,
                    text=ldgreen!55!black},
  lblr/.style     ={font=\scriptsize\sffamily, inner sep=1.5pt, text=ldred},
}

%% A red "no return path" marker, placed on a path with node[midway]{\xmark}.
\newcommand{\xmark}{\textcolor{ldred}{\sffamily\bfseries\large $\times$}}

%% ------------------------------------------------------------ title block
\titleformat{\section}{\normalfont\Large\bfseries\color{ldblue}}{\thesection}{0.7em}{}
\titleformat{\subsection}{\normalfont\large\bfseries\color{ldblue!85!black}}{\thesubsection}{0.6em}{}
\titleformat{\subsubsection}{\normalfont\normalsize\bfseries}{\thesubsubsection}{0.6em}{}
\setlist[itemize]{itemsep=2pt, topsep=3pt, parsep=0pt}
\setlist[enumerate]{itemsep=2pt, topsep=3pt, parsep=0pt}
\captionsetup{font=small, labelfont={bf,color=ldblue}}
.
\newcommand{\amp}{\prog{sender}}
\newcommand{\deamp}{\prog{receiver}}
%% The two ends of the link, named by ROLE rather than by security level.
%% A diode can be deployed in either orientation -- ingesting into a
%% protected domain, or exporting out of one -- so the mechanics here are
%% written in terms of which end sends and which receives. Each project's
%% "Orientation" section maps these onto its own security levels.
\newcommand{\srcside}{\textsc{source}}
\newcommand{\dstside}{\textsc{destination}}
\newcommand{\nm}[1]{#1\,nm}
\newcommand{\dbm}[1]{#1\,dBm}
\newcommand{\dbb}[1]{#1\,dB}

%% ------------------------------------------------------------ tikz styles
\tikzset{
  host/.style     ={draw, rounded corners=2pt, minimum width=27mm,
                    minimum height=12mm, align=center, font=\small,
                    fill=ldblue!8, draw=ldblue!70},
  mc/.style       ={draw, minimum width=25mm, minimum height=13mm,
                    align=center, font=\scriptsize, fill=gray!8,
                    draw=ldgray},
  netbox/.style   ={draw, dashed, rounded corners=3pt, draw=ldgray!70,
                    inner sep=5mm},
  splitter/.style ={draw, diamond, aspect=1.5, inner sep=1.5pt,
                    align=center, font=\scriptsize, fill=ldgreen!12,
                    draw=ldgreen},
  fiber/.style    ={-{Latex[length=2.4mm,width=1.8mm]}, thick,
                    draw=ldgreen!65!black},
  copper/.style   ={-, thick, draw=ldblue!70!black},
  dataflow/.style ={-{Latex[length=3.4mm,width=2.6mm]}, line width=1.5pt,
                    draw=ldblue},
  blocked/.style  ={-{Latex[length=2.4mm]}, thick, draw=ldred, dashed},
  lbl/.style      ={font=\scriptsize\sffamily, inner sep=1.5pt},
  lblg/.style     ={font=\scriptsize\sffamily, inner sep=1.5pt,
                    text=ldgreen!55!black},
  lblr/.style     ={font=\scriptsize\sffamily, inner sep=1.5pt, text=ldred},
}

%% A red "no return path" marker, placed on a path with node[midway]{\xmark}.
\newcommand{\xmark}{\textcolor{ldred}{\sffamily\bfseries\large $\times$}}

%% ------------------------------------------------------------ title block
\titleformat{\section}{\normalfont\Large\bfseries\color{ldblue}}{\thesection}{0.7em}{}
\titleformat{\subsection}{\normalfont\large\bfseries\color{ldblue!85!black}}{\thesubsection}{0.6em}{}
\titleformat{\subsubsection}{\normalfont\normalsize\bfseries}{\thesubsubsection}{0.6em}{}
\setlist[itemize]{itemsep=2pt, topsep=3pt, parsep=0pt}
\setlist[enumerate]{itemsep=2pt, topsep=3pt, parsep=0pt}
\captionsetup{font=small, labelfont={bf,color=ldblue}}
.
\newcommand{\amp}{\prog{sender}}
\newcommand{\deamp}{\prog{receiver}}
%% The two ends of the link, named by ROLE rather than by security level.
%% A diode can be deployed in either orientation -- ingesting into a
%% protected domain, or exporting out of one -- so the mechanics here are
%% written in terms of which end sends and which receives. Each project's
%% "Orientation" section maps these onto its own security levels.
\newcommand{\srcside}{\textsc{source}}
\newcommand{\dstside}{\textsc{destination}}
\newcommand{\nm}[1]{#1\,nm}
\newcommand{\dbm}[1]{#1\,dBm}
\newcommand{\dbb}[1]{#1\,dB}

%% ------------------------------------------------------------ tikz styles
\tikzset{
  host/.style     ={draw, rounded corners=2pt, minimum width=27mm,
                    minimum height=12mm, align=center, font=\small,
                    fill=ldblue!8, draw=ldblue!70},
  mc/.style       ={draw, minimum width=25mm, minimum height=13mm,
                    align=center, font=\scriptsize, fill=gray!8,
                    draw=ldgray},
  netbox/.style   ={draw, dashed, rounded corners=3pt, draw=ldgray!70,
                    inner sep=5mm},
  splitter/.style ={draw, diamond, aspect=1.5, inner sep=1.5pt,
                    align=center, font=\scriptsize, fill=ldgreen!12,
                    draw=ldgreen},
  fiber/.style    ={-{Latex[length=2.4mm,width=1.8mm]}, thick,
                    draw=ldgreen!65!black},
  copper/.style   ={-, thick, draw=ldblue!70!black},
  dataflow/.style ={-{Latex[length=3.4mm,width=2.6mm]}, line width=1.5pt,
                    draw=ldblue},
  blocked/.style  ={-{Latex[length=2.4mm]}, thick, draw=ldred, dashed},
  lbl/.style      ={font=\scriptsize\sffamily, inner sep=1.5pt},
  lblg/.style     ={font=\scriptsize\sffamily, inner sep=1.5pt,
                    text=ldgreen!55!black},
  lblr/.style     ={font=\scriptsize\sffamily, inner sep=1.5pt, text=ldred},
}

%% A red "no return path" marker, placed on a path with node[midway]{\xmark}.
\newcommand{\xmark}{\textcolor{ldred}{\sffamily\bfseries\large $\times$}}

%% ------------------------------------------------------------ title block
\titleformat{\section}{\normalfont\Large\bfseries\color{ldblue}}{\thesection}{0.7em}{}
\titleformat{\subsection}{\normalfont\large\bfseries\color{ldblue!85!black}}{\thesubsection}{0.6em}{}
\titleformat{\subsubsection}{\normalfont\normalsize\bfseries}{\thesubsubsection}{0.6em}{}
\setlist[itemize]{itemsep=2pt, topsep=3pt, parsep=0pt}
\setlist[enumerate]{itemsep=2pt, topsep=3pt, parsep=0pt}
\captionsetup{font=small, labelfont={bf,color=ldblue}}

\renewcommand{\amp}{\prog{od\_sender}}
\renewcommand{\deamp}{\prog{od\_receiver}}

\title{\vspace{-12mm}\textbf{Building a Data Diode: Quick Start}\\[2mm]
       \large A one-way OPC UA PubSub link in an afternoon}
\author{Alin-Adrian Anton \\
        \small \texttt{alin.anton@upt.ro} \\
        \small Politehnica University of Timi\c{s}oara}
\date{\small For the opc-diode project\\
      \small Diode build material shared across the sibling
      projects; canonical copy in log-diode}

\begin{document}
\maketitle
\thispagestyle{fancy}

\begin{abstract}
\noindent
This is the short road. It gets you from an empty bench to OPC UA PubSub NetworkMessages
crossing a genuinely unidirectional link, running the formally verified
\amp{} and \deamp{} tools at each end. Two routes are given: a
\textbf{software diode} for the bench --- fast, free, and \emph{not for
production} --- and a \textbf{hardware diode} built from three off-the-shelf
media converters, which is the real thing. Every claim about one-wayness is
paired with a test that proves it. For the reasoning behind the choices, the
optical power budgets, the two-converter/splitter variant and the threat
model, see the companion \emph{Data Diode Handbook}.
\end{abstract}

\vspace{2mm}
\tableofcontents
\vspace{4mm}

%% =====================================================================
\section{What you are building}

A data diode carries information one way between two networks and makes the
reverse physically impossible. Throughout this guide \srcside{} is the network
data leaves and \dstside{} the network it arrives in --- roles, not security
levels, because a diode is useful in both orientations: ingesting into a
protected domain, or exporting out of one.

Which end is the more trusted one depends on the deployment; the Handbook
\S2 states the orientation this project uses and what it protects.

\begin{figure}[h]
\centering
\resizebox{\textwidth}{!}{%
\begin{tikzpicture}[node distance=15mm]
  \node[host, minimum width=22mm] (src) {OPC UA\\publisher};
  \node[host, minimum width=22mm, right=of src] (amp) {\amp{}};
  \node[mc, right=of amp, minimum height=15mm, minimum width=20mm,
        fill=ldgreen!8, draw=ldgreen] (diode) {\textbf{DIODE}};
  \node[host, minimum width=22mm, right=of diode] (deamp) {\deamp{}};
  \node[host, minimum width=22mm, right=of deamp] (siem) {OPC UA\\subscribers};

  \draw[dataflow] (src)   -- (amp);
  \draw[dataflow] (amp)   -- node[lbl,above]{sealed} (diode);
  \draw[dataflow] (diode) -- node[lbl,above]{coded} (deamp);
  \draw[dataflow] (deamp) -- (siem);

  \draw[blocked] (siem.south) to[out=-140,in=-40]
        node[midway,fill=white,inner sep=1pt]{\xmark}
        node[lbl,below=3mm]{no return path, ever} (amp.south);

  \begin{scope}[on background layer]
    \node[netbox, inner sep=3mm, fit=(src)(amp),
          label={[lbl]above:SOURCE network}] {};
    \node[netbox, inner sep=3mm, fit=(deamp)(siem),
          label={[lbl]above:DESTINATION network}] {};
  \end{scope}
\end{tikzpicture}}
\caption{The shape of the system. Everything left of the diode is untrusted;
         everything right of it can never be reached from the left.}
\end{figure}

Because there is no return path there can be no acknowledgement, so there is
no TCP, no ARP resolution, no DNS and no retransmission. The transport must
be UDP, and reliability has to come from sending redundancy \emph{ahead of
time} rather than asking for a repeat. That is exactly what \amp{} and
\deamp{} do: forward error correction plus authenticated encryption, with the
parsing and reconstruction path machine-checked free of run-time errors.

\begin{notebox}[Which route should I take?]
Take the \textbf{software route} (\S\ref{sec:soft}) if you want to exercise
\amp{}/\deamp{} today, write tests, or demonstrate the pipeline. Take the
\textbf{hardware route} (\S\ref{sec:hard}) for anything that protects a real
network. The software route is a bench instrument, not a security control ---
see the warning at the head of \S\ref{sec:soft}.
\end{notebox}

%% =====================================================================
\section{Prerequisites}

\begin{itemize}
\item Two GNU/Linux hosts (or two VMs, or two network namespaces), each with a
      spare Ethernet interface for the diode link. In this guide the
      \srcside{} host is \cmd{192.168.10.1} and the \dstside{} host is
      \cmd{192.168.10.2}, both on \cmd{enp1s0}.
\item The project built: \cmd{./tools/build.sh} puts the sender and
      receiver binaries in \cmd{bin/}. This needs GNAT 14.2.0 and
      \prog{gprbuild}; see \file{tools/env.sh}.
\item Root on both hosts.
\item For the hardware route: three media converters, three simplex fiber
      strands, two Ethernet cables, and three power supplies. Bill of
      materials in \S\ref{sec:bom}.
\end{itemize}

Generate the shared key now --- you will need the same 64 hex characters on
both sides:

\begin{lstlisting}[style=shell]
openssl rand -hex 32
# e.g. 9f2c...4ab1   (32 bytes = 64 hex chars, ChaCha20-Poly1305)
\end{lstlisting}

\begin{warnbox}[The key cannot be exchanged over the diode]
There is no back channel, so there is no key agreement protocol available.
The key must be carried across by hand --- USB token, printed sheet, an
existing out-of-band management channel. Plan this before you rack anything.
\end{warnbox}

%% =====================================================================
\section{Route A --- the software diode (bench only)}
\label{sec:soft}

\begin{labbox}[Read this before using anything in this section]
A firewall-based diode enforces one-wayness with a \emph{mutable ruleset
inside a kernel that the attacker may be able to reach}. A root compromise,
a \cmd{nft flush ruleset}, a misapplied configuration-management run, or a
kernel bug restores the return path instantly and silently. It is a fine
instrument for developing against \amp{}/\deamp{}, for continuous
integration, and for teaching. \textbf{Do not deploy it as the boundary
between two security domains.} Use \S\ref{sec:hard}.
\end{labbox}

\subsection{Link-layer setup (both hosts)}

With no return path there is no ARP resolution, so the peer's MAC address has
to be configured by hand. This is the arrangement described in the paper
(\S2.2) and it is needed for the hardware diode too.

On the \srcside{} host, record the \dstside{} host's MAC:

\begin{lstlisting}[style=conf,caption={\file{/etc/ethers} on the SOURCE host}]
bb:bb:bb:bb:bb:bb 192.168.10.2
\end{lstlisting}

and on the \dstside{} host, record the \srcside{} host's MAC:

\begin{lstlisting}[style=conf,caption={\file{/etc/ethers} on the DESTINATION host}]
aa:aa:aa:aa:aa:aa 192.168.10.1
\end{lstlisting}

Then, on each host respectively:

\begin{lstlisting}[style=shell]
# SOURCE host
ip addr add 192.168.10.1/24 dev enp1s0
# DESTINATION host
ip addr add 192.168.10.2/24 dev enp1s0

# both hosts: load the static entries and make them permanent
arp -f /etc/ethers
ip neigh show dev enp1s0        # entries must read PERMANENT
\end{lstlisting}

Make it survive a reboot with a unit file (both hosts):

\begin{lstlisting}[style=conf,caption={\file{/etc/systemd/system/diode-link.service}}]
[Unit]
Description=Static ARP and addressing for the diode link
After=network-pre.target
Wants=network-pre.target

[Service]
Type=oneshot
RemainAfterExit=yes
ExecStart=/usr/sbin/arp -f /etc/ethers

[Install]
WantedBy=multi-user.target
\end{lstlisting}

\begin{lstlisting}[style=shell]
systemctl enable --now diode-link.service
\end{lstlisting}

Silence the interface so it cannot answer anything (both hosts):

\begin{lstlisting}[style=conf,caption={\file{/etc/sysctl.d/90-diode.conf}}]
net.ipv6.conf.enp1s0.disable_ipv6 = 1
net.ipv4.conf.enp1s0.arp_ignore   = 8
net.ipv4.conf.enp1s0.arp_announce = 2
net.ipv4.conf.enp1s0.accept_redirects = 0
net.ipv4.conf.enp1s0.send_redirects   = 0
net.ipv4.conf.enp1s0.accept_source_route = 0
\end{lstlisting}

\begin{lstlisting}[style=shell]
sysctl --system
\end{lstlisting}

\subsection{The ruleset}

The single most important rule is negative: \textbf{no connection tracking on
the diode interface}. A conntrack state entry, an \cmd{ESTABLISHED,RELATED}
accept, or a pf \cmd{keep state} \emph{is} a return path. So is an ICMP
port-unreachable and a TCP RST.

\paragraph{nftables (the modern default).} On the \srcside{} host, egress
only:

\begin{lstlisting}[style=nft,caption={\file{/etc/nftables.conf} --- SOURCE host}]
#!/usr/sbin/nft -f
flush ruleset

define DIODE_IF = "enp1s0"
define PEER     = 192.168.10.2
define PORT     = 6514

# Drop every frame arriving on the diode NIC, at the earliest possible
# point -- this catches ARP, LLDP and IPv6 too, which the inet family
# input hook would not see.
table netdev diodeguard {
    chain ingress {
        type filter hook ingress device enp1s0 priority -500; policy drop;
    }
}

table inet diode {
    chain input {
        type filter hook input priority filter; policy drop;
        iif lo accept
        iifname $DIODE_IF drop
    }
    chain output {
        type filter hook output priority filter; policy drop;
        oif lo accept
        oifname $DIODE_IF ip daddr $PEER udp dport $PORT accept
        oifname $DIODE_IF drop
    }
    chain forward {
        type filter hook forward priority filter; policy drop;
    }
}
\end{lstlisting}

On the \dstside{} host, ingress only:

\begin{lstlisting}[style=nft,caption={\file{/etc/nftables.conf} --- DESTINATION host}]
#!/usr/sbin/nft -f
flush ruleset

define DIODE_IF = "enp1s0"
define PEER     = 192.168.10.1
define PORT     = 6514

table inet diode {
    chain input {
        type filter hook input priority filter; policy drop;
        iif lo accept
        iifname $DIODE_IF ip saddr $PEER udp dport $PORT accept
        iifname $DIODE_IF drop
    }
    chain output {
        type filter hook output priority filter; policy drop;
        oif lo accept
        # Nothing may ever leave by the diode NIC.
        oifname $DIODE_IF drop
    }
    chain forward {
        type filter hook forward priority filter; policy drop;
    }
}
\end{lstlisting}

\begin{lstlisting}[style=shell]
systemctl enable --now nftables
nft list ruleset          # read it back and check it is what you wrote
\end{lstlisting}

\paragraph{iptables and pf.} Equivalent rulesets for legacy \prog{iptables},
OpenBSD \prog{pf} and FreeBSD \prog{pf} --- including the
\cmd{set block-policy drop} subtlety that stops pf from answering with RSTs
--- are in the Handbook, \S4.

\subsection{Route A$'$ --- a unidirectional layer-2 tunnel}
\label{sec:tap}

There is a better software option than filtering at layer 3. Bridge a TAP
interface to each of two \emph{real} Ethernet segments and forward frames
between the taps with a program that has no reverse code path. You get the
same Ethernet semantics as the optical diode of \S\ref{sec:hard} --- a
one-way layer-2 tunnel --- with software in place of the missing transmitter.

\begin{lstlisting}[style=shell,caption={On the forwarding host}]
ip tuntap add dev tapdlow  mode tap
ip tuntap add dev tapdhigh mode tap

# One bridge per side.  STP off -- BPDUs are bidirectional signalling.
ip link add br-low  type bridge stp_state 0
ip link add br-high type bridge stp_state 0
ip link set eth0     master br-low     # NIC facing the SOURCE segment
ip link set tapdlow  master br-low
ip link set eth1     master br-high    # NIC facing the DESTINATION segment
ip link set tapdhigh master br-high
for i in eth0 eth1 tapdlow tapdhigh br-low br-high; do ip link set "$i" up; done

./tapdiode tapdlow tapdhigh            # forwards low -> high, and only that
\end{lstlisting}

The forwarder is about a hundred lines of C; the whole program, with notes on
its design, is in the Handbook (\S4.7). Its core is the part worth seeing
here: only the input descriptor is ever polled, and \cmd{read()} is never
called on the output descriptor, so the reverse direction is absent rather
than forbidden.

\begin{lstlisting}[style=c,caption={The loop, in essence}]
pfd.fd = in_fd;           /* the ONLY descriptor ever polled */
pfd.events = POLLIN;

while (!stop_requested) {
    if (poll(&pfd, 1, -1) < 0) { if (errno == EINTR) continue; break; }

    n = read(in_fd, frame, sizeof frame);        /* only ever in_fd  */
    if (n <= 0) { if (errno == EINTR) continue; break; }

    if (write(out_fd, frame, (size_t) n) != n)   /* only ever out_fd */
        lost++;                                  /* loss, never feedback */
}
\end{lstlisting}

\begin{warnbox}[Two bridges, never one]
Everything rests on \cmd{br-low} and \cmd{br-high} being separate. Putting
\cmd{eth0} and \cmd{eth1} in the \emph{same} bridge gives an ordinary
bidirectional bridge and destroys the property --- and it is one command.
Assert it afterwards: \cmd{bridge link show} must never list both NICs under
one master. This is still software, so \S\ref{sec:hard} remains the answer for
a real boundary.
\end{warnbox}

Verify at layer 2 --- \cmd{ping} proves nothing about frames the IP stack
never generates. Sniff \cmd{'ether proto 0x88b5'} on the \dstside{} segment
with \cmd{tcpdump} while injecting one tagged frame on \srcside{}, then swap
the two and confirm nothing arrives (Handbook \S4.7 has both commands).

%% =====================================================================
\section{Route B --- the hardware diode}
\label{sec:hard}

This is the three-converter arrangement. It needs no fiber splitter, works
with every model tested, and is the one to build if you are unsure.

\subsection{Why three converters}

A media converter bridges copper Ethernet and fiber. It has two \emph{separate}
fiber ports, TX and RX. Left with nothing on its RX port, a converter sees no
carrier, declares the fiber link down, and stops forwarding --- so you cannot
simply omit the return fiber. The trick is to give the sender light that
carries \emph{no data from the other side}. A third converter, its copper port
deliberately left unplugged, supplies exactly that: carrier with nothing on it.

\begin{figure}[h]
\centering
\resizebox{\textwidth}{!}{%
\begin{tikzpicture}
  \node[host, minimum width=22mm] (low) {SOURCE host\\\scriptsize 192.168.10.1};
  \node[mc, minimum width=22mm, right=14mm of low]
       (A) {\textbf{MC-A} (input)\\\scriptsize TX \quad RX};
  \node[mc, minimum width=22mm, right=24mm of A]
       (B) {\textbf{MC-B} (output)\\\scriptsize TX \quad RX};
  \node[host, minimum width=22mm, right=14mm of B]
       (high) {DESTINATION host\\\scriptsize 192.168.10.2};
  \node[mc, minimum width=30mm, below=18mm of $(A)!0.5!(B)$,
        fill=ldamber!10, draw=ldamber]
       (C) {\textbf{MC-C} (emulator)\\\scriptsize copper port UNUSED};

  \draw[copper] (low) -- node[lbl,above]{Cat5e} (A);
  \draw[copper] (B) -- node[lbl,above]{Cat5e} (high);

  \draw[fiber] (A.north east) to[out=40,in=140]
        node[lblg,above]{\textbf{1.} TX $\rightarrow$ RX \emph{(data path)}} (B.north west);
  \draw[fiber] (B.south) to[out=-90,in=0]
        node[lblg,right,pos=0.35]{\textbf{2.} TX $\rightarrow$ RX} (C.east);
  \draw[fiber] (C.west) to[out=180,in=-90]
        node[lblg,left,pos=0.65]{\textbf{3.} TX $\rightarrow$ RX} (A.south);

  \node[lblr, below=1mm of C] {copper unplugged $\Rightarrow$ MC-C carries carrier, never data};
\end{tikzpicture}}
\caption{Three-converter diode. Data travels only along link~1. Links~2 and~3
         exist solely so MC-A sees carrier on its RX port; because MC-C's
         copper port is unplugged, no data can traverse them.}
\label{fig:triangle}
\end{figure}

\begin{warnbox}[The one wire that must stay unplugged]
MC-C's \emph{copper} port must be left disconnected. If anyone patches it into
the \srcside{} network, links 2 and 3 become a complete return path from
\dstside{} to \srcside{} and the diode is gone. Label that port, cover it, and
put it on the commissioning checklist.
\end{warnbox}

\subsection{Bill of materials}
\label{sec:bom}

Any dual-fiber media converter with separate TX and RX ports will do. These
three are the ones this project has been built and tested with, and they cover
fast Ethernet, gigabit multimode and gigabit single-mode.

\begin{table}[h]
\centering\small
\begin{tabularx}{\textwidth}{@{}p{25mm}llllX@{}}
\toprule
\textbf{Model} & \textbf{Speed} & \textbf{Fiber} & \textbf{$\lambda$} &
\textbf{Conn.} & \textbf{Notes} \\
\midrule
TP-Link\newline MC100CM & 100\,Mb/s & multi-mode  & \nm{1310} & SC &
  Fast Ethernet, up to 2\,km. Works in both topologies; \nm{1310} splitters
  are catalogue parts. \\
TP-Link\newline MC200CM & 1\,Gb/s   & multi-mode  & \nm{850}  & SC &
  Gigabit, up to 550\,m on OM2. Works in both topologies --- but an \nm{850}
  splitter is a special order with a long lead time. \\
TP-Link\newline MC210CS & 1\,Gb/s   & single-mode & \nm{1310} & SC &
  Gigabit. TP-Link's own sources give both 15\,km and 20\,km --- check the
  revision you buy. Cheap, commodity \nm{1310} splitters. \\
\bottomrule
\end{tabularx}
\caption{Distances are vendor figures for a normal link; a diode's fibers are
         a metre long, so see the note on attenuators below. Use three
         converters of the \emph{same} model --- mixing speeds or wavelengths
         across the triangle will not link up.}
\end{table}

The triangle uses \textbf{three simplex fiber strands} --- links 1, 2 and 3 in
\S\ref{sec:hard}. Duplex patch cords carry two strands each, so two duplex
cords cover it with one strand spare. Match the connector and fiber type to
the converter, and budget one power supply per converter (9\,V DC,
centre-positive barrel). See Handbook \S5.6 for power, UPS and rack notes.

\begin{tipbox}[Attenuators: only if you use long-range optics]
Three \emph{matched} converters need no attenuation, however short the fiber:
within an optical class the maximum launch power equals the maximum receive
power, so a compliant transmitter cannot overload a compliant receiver of the
same class. If instead you use a long-range converter or SFP --- built to drive
tens of kilometres --- across the metre of fiber a diode needs, it can exceed
the receiver's maximum input power, which shows up as CRC errors and dropped
frames. Fixed attenuators (\dbb{3}, \dbb{5}, \dbb{10}) are the fix; size them
from the datasheet figures. Handbook \S5.3.
\end{tipbox}

\subsection{Wiring}

\begin{table}[h]
\centering\small
\begin{tabularx}{\textwidth}{@{}llX@{}}
\toprule
\textbf{From} & \textbf{To} & \textbf{Why} \\
\midrule
SOURCE host NIC        & MC-A copper       & data enters the diode \\
\textbf{MC-A fiber TX} & \textbf{MC-B fiber RX} & \textbf{the one-way data path} \\
MC-B copper         & DESTINATION host NIC     & data leaves the diode \\
MC-B fiber TX       & MC-C fiber RX     & keeps MC-C's transceiver up \\
MC-C fiber TX       & MC-A fiber RX     & gives MC-A carrier so it holds link \\
MC-C copper         & \textbf{nothing}  & breaks the loop; \emph{leave unplugged} \\
MC-A fiber RX       & from MC-C only    & must never come from MC-B \\
\bottomrule
\end{tabularx}
\end{table}

Power all three converters, then confirm each shows both a copper link LED and
a fiber link LED (except MC-C, whose copper LED stays dark --- that is
correct). Now apply the \emph{same} link-layer setup as in \S\ref{sec:soft}:
addresses, \file{/etc/ethers}, \cmd{arp -f}, and the sysctl file. The firewall
rules are optional here --- the optics already enforce the direction --- but
applying them anyway is good defence in depth.

\begin{tipbox}
If MC-C goes completely dark with its copper port unplugged, your converters
implement link-fault pass-through. Give MC-C a copper link that goes nowhere:
an isolated switch port on a dead VLAN, or a loopback plug. Never give it a
port on the \srcside{} network.
\end{tipbox}

%% =====================================================================
\section{Prove it is one-way}
\label{sec:prove}

Do not skip this. Run all five, and record the output --- this is your
commissioning evidence.

\begin{enumerate}
\item \textbf{Physical trace.} Follow every fiber by hand against the wiring
      table. Confirm MC-C's copper port is unplugged.

\item \textbf{Carrier independence.} Unplug MC-B's copper cable from the
      \dstside{} host. MC-A must \emph{keep} its fiber link LED. If MC-A goes
      down, its carrier is coming from the far end and you have built a
      bidirectional link.

\item \textbf{Reverse probe.} On the \srcside{} host, start a capture that
      accepts everything:
\begin{lstlisting}[style=shell]
# SOURCE host -- must stay silent for the whole test
tcpdump -ni enp1s0 -e -vv
\end{lstlisting}
      From the \dstside{} host, try hard to get a packet back:
\begin{lstlisting}[style=shell]
# DESTINATION host
ping    -c 5           192.168.10.1
arping  -c 5 -I enp1s0 192.168.10.1
nc -u -w2              192.168.10.1 9999 <<< "probe"
hping3 -S -c 5 -p 22   192.168.10.1        # TCP SYN
\end{lstlisting}
      The \srcside{} capture must show \textbf{nothing} from the \dstside{}
      host. Every command above must fail or time out.

\item \textbf{TCP impossibility.} From the \dstside{} host,
      \cmd{nc -w 5 192.168.10.1 22} must always time out. A three-way
      handshake cannot complete across a diode; if it does, you have a
      bidirectional link.

\item \textbf{Forward path works.} From the \srcside{} host, send a datagram
      and see it arrive:
\begin{lstlisting}[style=shell]
# DESTINATION host
tcpdump -ni enp1s0 udp port 6514
# SOURCE host
echo hello | nc -u -w1 192.168.10.2 6514
\end{lstlisting}
\end{enumerate}

\begin{notebox}[There is no remote diagnostic]
Unmanaged converters cannot tell you they have failed, and nothing can report
back across the diode. Schedule a physical inspection --- the Handbook \S6.5
has a checklist.
\end{notebox}

%% =====================================================================
\section{Run the tools}

Start the receiving side first, so nothing is lost.

\begin{lstlisting}[style=shell,caption={DESTINATION side}]
./bin/od_receiver 9702 127.0.0.1 9703 --key 9f2c...4ab1
\end{lstlisting}

\begin{lstlisting}[style=shell,caption={SOURCE side}]
./bin/od_sender 9701 192.168.10.2 9702 --parity 3 --interleave 8 \
                --pace-us 200 --key 9f2c...4ab1
\end{lstlisting}

The publisher points its PubSub output at the sender's input port;
subscribers listen on the receiver's output port.

The Handbook \S7 explains the redundancy settings, and the project
README is the authority on the full command line.

\subsection{As services}

\begin{lstlisting}[style=conf,caption={\file{/etc/systemd/system/od-sender.service}}]
[Unit]
Description=opc-diode sender (SOURCE side)
After=network-online.target diode-link.service
Wants=network-online.target

[Service]
Type=simple
ExecStart=/opt/opc-diode/bin/od_sender 9701 192.168.10.2 9702 \\
          --parity 3 --key ${DIODE_KEY}
EnvironmentFile=/etc/opc-diode/key.env
Restart=always
RestartSec=2
DynamicUser=yes
NoNewPrivileges=yes
ProtectSystem=strict
ProtectHome=yes
PrivateTmp=yes

[Install]
WantedBy=multi-user.target
\end{lstlisting}

Keep the key (and seed, where one is used) in a mode-\cmd{0400}
\cmd{EnvironmentFile} rather than on the command line, so it does not
appear in \cmd{ps}.

%% =====================================================================
\section{Where to go next}

\begin{itemize}
\item \textbf{The two-converter build.} An ordinary $1\times2$ 50:50 splitter
      replaces the third converter, and works at every wavelength
      here.\footnote{The \emph{Sensors} paper states that this approach ``did
      not work'' for the MC200CM, and reports the three-converter arrangement
      being used instead (Anton et al., \S2.2). That wording describes the
      parts available at the time rather than a limitation of the method: no
      \nm{850} multimode splitter could be obtained within the timeframe of
      that work. The one eventually used had to be made to special order and
      took months to arrive, and with it in hand the two-converter topology
      works on the MC200CM exactly as it does at \nm{1310}.} What differs is
      procurement --- \nm{1310} splitters are catalogue stock, \nm{850} ones
      are a special order. Handbook \S5.2--5.3.
\item \textbf{Embedded and IoT.} Between two microcontrollers the diode is
      three wires: SPI with the MISO conductor omitted, optionally through
      optocouplers so each line is unidirectional in hardware. Handbook
      \S5.8 --- including why the \emph{sender} must be the bus master.
\item \textbf{Threat model.} What a diode does and does not protect against,
      including the insider and covert-channel discussion --- Handbook \S2.
\item \textbf{iptables and pf.} Full equivalents of the nftables ruleset,
      Handbook \S4.
\item \textbf{The layer-2 tunnel in full.} The complete \file{tapdiode.c},
      its design notes, the bridge and \prog{sysctl} hardening, and how to
      verify it at layer 2 --- Handbook \S4.7.
\item \textbf{Why the daemons are proved.} The parser on the \dstside{} side
      chews bytes chosen by whoever compromised the \srcside{} network, on a
      link with no way to ask for a retry --- Handbook \S7 and
      \file{docs/ASSURANCE.md}.
\item \textbf{Scaling.} Several diodes in parallel, round-robin ---
      Handbook \S5.7.
\end{itemize}

\vfill
\begin{center}\small
\textcolor{ldgray}{%
Made \copyleft{} libre (free as in freedom) at
Politehnica University of Timi\c{s}oara.\\
Copyright \copyright{} 2026 Alin-Adrian Anton.\\
This document is licensed under the Creative Commons Attribution 4.0
International Licence (CC~BY~4.0) or any later version.\\
\url{https://creativecommons.org/licenses/by/4.0/}}
\end{center}

\end{document}
